\section{Introduction}

In the standard $\Lambda$CDM model the statistical isotropy of the universe is considered as one of the main foundations. This idea is supported by the simplest slow-roll inflationary scenarios, where the primordial fluctuations, the seeds for late universe structures, are a single realization of a statistically isotropic Gaussian random field \cite{2020A&A...641A..10P}. These properties can be accurately tested trough the Cosmic Microwave Background (CMB) anisotropies.

Observations of the CMB by the Planck satellite are consistent with the Gaussian, almost scale-invariant, and adiabatic nature of the primordial perturbations predicted by the slow-roll inflation, providing a very accurate confirmation of the base $\Lambda CDM$ model \cite{Planck:2018nkj, 2020A&A...641A...5P, 2020A&A...641A...6P, Planck:2019evm}. However, some unexpected features on large angular scales have been found in the temperature fluctuation field. These anomalies were already observed soon after the first release of the NASA \textit{Wilkinson Microwave Anisotropy Probe} (WMAP) satellite data \cite{2003ApJS..148....1B, 2007ApJS..170..288H, 2009ApJS..180..225H, 2011ApJS..192...17B, 2013ApJS..208...19H}, and they show tensions with the standard model. Some of them suggest a potential violation of the cosmological principle, while others are related to non-Gaussian imprints. Among these features are the power asymmetry \cite{2004ApJ...605...14E, 2004MNRAS.354..641H}, quadrupole-octopole alignment\cite{2004PhRvD..69f3516D}, the cold spot \cite{2004ApJ...609...22V, 2010AdAst2010E..77V}, the anomalously low variance \citep{Monteserin08} or the lack of power at large scales \cite{2006A&A...454..409B}. In \cite{2023CQGra..40i4001K} a summary of some of the anomalies observed in the CMB can be found.


In the present work, we are interested in the "hemispherical power asymmetry" (HPA). This anomalous feature has been found with different methods at the level of 3$\sigma$-3.5$\sigma$ depending on the estimator and the angular scales used for the analysis \cite{2007ApJ...660L..81E, 2009ApJ...704.1448H, 2009ApJ...699..985H, 2014A&A...571A..23P, Akrami:2014eta, 2015MNRAS.446.4232A, 2016A&A...594A..16P, Planck:2019evm}. The asymmetry persists in Planck data, ruling out systematics as a plausible explanation, as it has been observed at a similar level by two independent missions.

The HPA has been modeled as a dipole modulation of the primordial fluctuations \cite{2005PhRvD..72j3002G, 2007ApJ...656..636G} as
%
\begin{equation}\label{DM}
    \Delta{\tilde{T}}(\hat{n}) = (1+A_{DM}\hat{n}\cdot\hat{p})\Delta{T}(\hat{n})
\end{equation}
%
where $\Delta{\tilde{T}}$ and $\Delta{T}$ are the modulated and isotropic CMB fluctuation fields, respectively, and $A_{DM}$ and $\hat{p}$ the amplitude and direction of the modulation.

Some analyses have 
performed a direct fitting
to this model \cite{2009ApJ...699..985H, 2007ApJ...660L..81E, 2014A&A...571A..23P, 2017PhRvD..95f3011Z}. They found amplitudes close to 7\% 
at a confidence level
around 3$\sigma$. Equivalent results have been obtained following methods working in 
spherical harmonics space \cite{2015PhRvD..91b3515R, 2009PhRvD..80f3004H, 2016JCAP...01..046G, 2016A&A...594A..16P}. These estimators exploit the expected correlations between $\ell$ and $\ell\pm 1$ induced by the possible dipolar modulation. The HPA has also been formulated in terms of Bipolar spherical harmonics (BipoSH) \cite{2016A&A...594A..16P, 2023arXiv230104539D} technique, which works in the harmonic space and has the capability of tracing 
statistical isotropy violations.

It has also been observed that the dipolar modulation is not scale-independent in the sense that the amplitude is much lower for small angular scales and apparently vanishes for multipoles above six hundred \cite{2009ApJ...704.1448H, Akrami:2014eta, Marcos-Caballero:2019jqj}. In particular, in most analyses, the statistical significance decreases if the full range of scales is taken into account.
In any case, the dipolar power asymmetry remains an intriguing anomaly present on super-Hubble scales at the recombination time, suggesting a possible physical origin related to the inflationary mechanism \cite{2016JCAP...01..046G, 2016A&A...594A..16P}.

Although most of the analyses have assummed a dipolar modulation as the fitting model, it is important to remark that this is not necessarily the underlying model of the HPA. There have been some attempts to propose physical models that could explain the origin of the HPA and reconcile observations with the standard cosmological model \cite{2008PhRvD..78l3520E, 2009PhRvD..80b3526D, 2010PhRvD..81h3501C, 2011JHEP...02..061K, 2013JCAP...08..007L, 2016MNRAS.460.1577K, 2017PhRvD..95f3011Z, 2022arXiv220903928S}. Some of them suggest a superhorizon perturbation as a remnant of a pre-inflationary phase where inflation is given in a multi-field scenario, including curvaton-type models \cite{2008PhRvD..78l3520E}. This type of models predicts non-Gaussian imprints in the CMB maps, which can be used to constrain these scenarios. However, according to the latest Planck inflationary constraints, 
they are not statistically favored compared to $\Lambda$CDM \cite{2020A&A...641A..10P}. A non-cosmological explanation has been recently proposed related to a new foreground associated to nearby galaxies \cite{Hansen:2023gra}.

In this paper, we perform an analysis similar to the one carried out by the Planck collaboration in \cite{Planck:2019evm} using the latest Planck Release 4 (PR4) data set, also known as NPIPE, in both intensity and polarization E-mode maps. In particular, we use a pixel space \emph{local-variance} method introduced in \cite{Akrami:2014eta}. We study the presence of the HPA in the intensity clean CMB maps 
obtained from PR4 using the Sevem \cite{2012MNRAS.420.2162F} and Commander \cite{2008ApJ...676...10E} componet separation methods. In polarization only Sevem data is used.

Exploring the asymmetry in polarization data is very important to answer the question of whether the HPA is related to new physics or a statistical fluke, as the large-scale intensity data on its own is not able to claim a strong positive detection and is already cosmic-variance limited. Furthermore, as discussed in \cite{2020MNRAS.492.3994G}, different proposed models for the HPA in intensity will also predict a similar asymmetry in CMB polarization data.
There have been some attempts to estimate the significance level of the HPA in the Planck E-mode map \cite{2017arXiv171000580A, 2020MNRAS.492.3994G, Planck:2019evm}. Using the PR3 data set, Planck collaboration found a large dispersion in the p-value between the four component separation methods, reflecting the presence of different amounts of residuals in the corresponding maps. Actually, Sevem and Commander returned a p-value below 1\%, while NILC and SMICA showed a p-value closer to 5\%-6\%.

The structure of the paper is as follows. In section \ref{DatSim} we describe the data and simulations used in the analysis. The methodology is outlined in section \ref{Met}, where we describe in detail the inpainting approach, the local-variance estimator, and the procedure followed for generating the masks. The main results are shown in Section \ref{Res}. Finally, a summary of the main conclusions from the analysis is given in section \ref{Con}.
